\documentclass[12pt]{article}
\usepackage[utf8]{inputenc}
\usepackage[spanish]{babel}
\usepackage{amsmath}
\usepackage{geometry}
\usepackage{booktabs}

\geometry{a4paper, margin=1in}

\title{Propuesta de Proyecto: Optimización de Logística y Distribución}
\author{Investigación de Operaciones}
\date{\today}

\begin{document}

\maketitle

\section{Contexto del Problema}
En el entorno competitivo de la logística urbana, la eficiencia en la ``última milla'' es crucial. La empresa \textbf{Logística Distrital S.A.} debe distribuir suministros tecnológicos en Bogotá (Usaquén, Chapinero y Kennedy). Se requiere un modelo que optimice la asignación de flota vehicular para cumplir con las entregas mínimas minimizando los costos operativos.

\section{Planteamiento del Problema}
Se cuenta con una flota de $10$ camiones. El objetivo es determinar la asignación óptima de vehículos ($x_i$) para tres zonas con diferentes requerimientos:

\begin{table}[h]
\centering
\begin{tabular}{@{}ccccc@{}}
\toprule
\textbf{Zona ($i$)} & \textbf{Nombre} & \textbf{Tiempo ($t_i$)} & \textbf{Paquetes ($p_i$)} & \textbf{Demanda ($d_i$)} \\ \midrule
1 & Usaquén & 60 min & 15 & 45 \\
2 & Chapinero & 45 min & 10 & 30 \\
3 & Kennedy & 90 min & 20 & 60 \\ \bottomrule
\end{tabular}
\end{table}

\section{Modelo Matemático: Método de la Gran M}
\subsection{Variables de Decisión}
\begin{itemize}
    \item $x_i$: Vehículos asignados a la zona $i$.
    \item $a_i$: Variables artificiales de penalización.
\end{itemize}

\subsection{Función Objetivo}
\[ \text{Min } Z = 60x_1 + 45x_2 + 90x_3 + M(a_1 + a_2 + a_3) \]

\subsection{Restricciones}
\begin{align}
    x_1 + x_2 + x_3 &\le 10 \quad \text{(Flota)} \\
    15x_1 + a_1 &\ge 45 \quad \text{(Demanda 1)} \\
    10x_2 + a_2 &\ge 30 \quad \text{(Demanda 2)} \\
    20x_3 + a_3 &\ge 60 \quad \text{(Demanda 3)} \\
    x_i, a_i &\ge 0
\end{align}

\section{Modelo Matemático: Programación Dinámica}
\begin{itemize}
    \item \textbf{Etapa ($n$):} Zonas 1, 2 y 3.
    \item \textbf{Estado ($s_n$):} Vehículos disponibles para la etapa $n$.
    \item \textbf{Decisión ($x_n$):} Vehículos asignados a la zona $n$.
\end{itemize}

\subsection{Ecuación de Recurrencia}
\[ f_n(s_n) = \max_{0 \le x_n \le s_n} \{ p_n(x_n) + f_{n+1}(s_n - x_n) \} \]

\end{document}
